\section{Related Work}

\paragraph{Low-rank tensor models and Tucker backbones.}
CP, Tucker, tensor train, and tensor ring decompositions provide compact
representations for multiway data through different rank notions and factor
topologies \citep{Kolda2009Tensor,Sidiropoulos2017TSP,
DeLathauwer2000MLSVD,oseledets2011,zhao2016tr}. Recent completion work has
improved scalability, sampling guarantees, and visual-data priors while still
mostly changing the tensor rank model, tensor-network topology, or low-rank
regularizer \citep{LiuMoitra2020TensorCompletion,Wang2021CoupledTransform,
Zheng2021FCTN}. Tucker is the closest baseline for NTD-PL because it uses the
same core-factor multilinear backbone. NTD-PL does not primarily compete by
changing the rank tuple, replacing the factor topology, or moving to another
tensor-network format. Instead, it keeps the Tucker backbone and multilinear
ranks fixed and changes the measurement map on top of the Tucker latent tensor.

\paragraph{Nonlinear tensor factorization and neural/kernel decoders.}
Nonlinear tensor models include generalized likelihood and link formulations,
Bayesian and kernelized tensor factorizations, Gaussian-process tensor models,
and neural tensor decoders
\citep{hong2020generalized,schein2016bayesian,xu2011infinite,
zhe2016distributed,liu2017deepcp,pan2020streaming,
saragadam2024deeptensor,varolgunes2023nmtucker,ahn2024neural,
Luo2024LRTFR}. Related imaging methods also learn linear or nonlinear
transforms under low-rank tensor priors
\citep{Luo2022SelfSupervisedTNN,Wang2024CoNoT,xu2024stacked}. These methods
can be more flexible than a scalar polynomial link, but this flexibility often
comes with heavier inference, additional hyperparameters, or a less transparent
finite-dimensional relation to Tucker. NTD-PL is complementary: it deliberately
uses a low-dimensional scalar response rather than a general kernel,
Gaussian-process, neural decoder, or learned transform, so that the induced
effect can be analyzed relative to a fixed Tucker backbone.

\paragraph{Polynomial feature lifts and high-order interactions.}
Polynomial and Volterra-style expansions are classical tools for representing
nonlinear responses \citep{Volterra1959Theory,Boyd1985FadingMemory}. Related
matrix and tensor methods use polynomial feature maps, quadratic
factorizations, or explicit high-order interactions to enlarge a linear model
\citep{zhai2024quadratic,Ruegamer2024AISTATS,liu2015model,deng2016storm}.
Such constructions can make nonlinear interactions explicit, but untied
high-order parameterizations can increase parameter dimension and contraction
cost. NTD-PL does not introduce an independent high-order tensor parameter. Its
polynomial powers are powers of one learned Tucker latent tensor, inducing
parameter-tied high-order interactions and Veronese-style backbone-tied rank
lifts.

\paragraph{Positioning and scope.}
NTD-PL is not meant to replace broad nonlinear tensor decoders or higher-rank Tucker models.
It targets a specific regime: the Tucker backbone remains useful, but matched-rank Tucker leaves
a residual that is coherent with a nonlinear response. In this regime, a polynomial link provides
a compact rank lift tied to the same latent tensor, rather than an untied high-rank correction or a
general neural/kernel decoder.
